% Lista de acrónimos 
% Para llamarlos usar \acrshort
\newacronym{aes}{AES}{Advanced Encryption Standard}
\newacronym{api}{API}{Application Programming Interface}
\newacronym{captcha}{CAPTCHA}{Completely Automated Public Turing test to tell Computers and Humans Apart}
\newacronym{cbc}{CBC}{Cipher Block Chaining}
\newacronym{cicd}{CI/CD}{Continuous Integration / Continuous Delivery}
\newacronym{cpu}{CPU}{Central Processing Unit}
\newacronym{crud}{CRUD}{Create, Read, Update, Delete}
\newacronym{csrf}{CSRF}{Cross-Site Request Forgery}
\newacronym{css}{CSS}{Cascading Style Sheets}
\newacronym{db}{DB}{Database}
\newacronym{vdom}{VDOM}{Virtual Document Object Model}
\newacronym{ecb}{ECB}{Electronic Codebook}
\newacronym{ecc}{ECC}{Elliptic Curve Cryptography}
\newacronym{javaee}{Java EE}{Java Enterprise Edition}
\newacronym{gb}{GB}{Gigabyte}
\newacronym{gcm}{GCM}{Galois/Counter Mode}
\newacronym{get}{GET}{Hypertext Transfer Protocol GET Method}
\newacronym{html}{HTML}{HyperText Markup Language}
\newacronym{http}{HTTP}{HyperText Transfer Protocol}
\newacronym{ide}{IDE}{Integrated Development Environment}
\newacronym{idea}{IDEA}{International Data Encryption Algorithm}
\newacronym{jdk}{JDK}{Java Development Kit}
\newacronym{jpa}{JPA}{Java Persistence API}
\newacronym{jpql}{JPQL}{Java Persistence Query Language}
\newacronym{json}{JSON}{JavaScript Object Notation}
\newacronym{jsx}{JSX}{JavaScript XML}
\newacronym{mvc}{MVC}{Model-View-Controller}
\newacronym{post}{POST}{Hypertext Transfer Protocol POST Method}
\newacronym{put}{PUT}{Hypertext Transfer Protocol PUT Method}
\newacronym{ram}{RAM}{Random Access Memory}
\newacronym{rest}{REST}{Representational State Transfer}
\newacronym{rsa}{RSA}{Rivest–Shamir–Adleman (Cryptosystem)}
\newacronym{sid}{SID}{Sistema de Impresión de Documentos}
\newacronym{smtp}{SMTP}{Simple Mail Transfer Protocol}
\newacronym{spa}{SPA}{Single Page Application}
\newacronym{sql}{SQL}{Structured Query Language}
\newacronym{tls}{TLS}{Transport Layer Security}
\newacronym{uml}{UML}{Unified Modeling Language}
\newacronym{url}{URL}{Uniform Resource Locator}
\newacronym{xml}{XML}{eXtensible Markup Language}


\newglossaryentry{policies}{
    name=policies,
    description={son las reglas y modelos que determinan cómo el asistente virtual decide cuál será la próxima acción en una conversación, basándose en el contexto y el historial del diálogo.}
}

\newglossaryentry{pipeline}{
    name=pipeline,
    description={secuencia de componentes y pasos configurados para procesar datos, extraer características, y entrenar modelos que permiten al asistente virtual comprender e interpretar el lenguaje natural. En caso de Rasa se encuentra en config.yaml}
}

\newglossaryentry{metodologia agil}{
    name=metodología ágil,
    description={
        En el desarrollo de software, las prácticas ágiles consisten en la mejora de requisitos, investigación y soluciones mediante el esfuerzo colaborativo de equipos autoorganizados y multifuncionales junto con sus clientes/usuarios finales.
    }
}

